\begin{claim}{C.18}{Refuted}{The complete interval hypotheses do not imply exact coproducts}\label{C18}
The manuscript's (a)--(d), existence of all small limits and colimits, and even cocontinuity of the tensor product in both variables do not imply AB4 or AB5. The counterexample below still satisfies the asserted underlying localization equivalence. Thus it refutes a proposed route to the general theorem, not that theorem.
\end{claim}
\begin{proof}
We work over $k=\mathbb Z$. Let $\mathcal T$ be the category of torsion abelian groups and put $\mathcal N=\mathcal T^{\mathrm{op}}$. Write $tA$ for the torsion subgroup of an abelian group $A$. The category $\mathcal T$ is abelian, with ordinary kernels and cokernels, ordinary direct sums, and products $t(\prod_iA_i)$. This product formula is also recorded in \citet[Introduction]{Rickard2020C}. Consequently $\mathcal T$ and $\mathcal N$ have all small limits and colimits.

The category $\mathcal T$ has enough injectives. For an explicit verification, characters into $\mathbb Q/\mathbb Z$ separate the points of a torsion group $A$: a nonzero character on its cyclic subgroup extends to $A$, since $\mathbb Q/\mathbb Z$ is divisible and hence injective. The resulting embedding of $A$ into a product of copies of $\mathbb Q/\mathbb Z$ lands in its torsion subgroup. That subgroup is divisible: if $nx=y$ is solved in the divisible product and $ry=0$, then $nrx=0$, so the solution $x$ is torsion. It is therefore injective even in all abelian groups, and hence in $\mathcal T$. We use here the divisible--injective characterization, with its proof opened in \citet[Lemma 15.55.1, tag 01D6]{StacksInjectiveAbC}. Thus $\mathcal N$ has enough projectives.

The category $\mathcal N$ has its usual scalar copowers. Explicitly, for an abelian group $U$ and an object $A$ of $\mathcal N$, define
\[
 U\cdot A=t\operatorname{Hom}_{\mathbb Z}(U,A),
\]
regarded as an object of $\mathcal N$. The opposite variance turns this into a covariant functor in $U$ and $A$. Indeed
\[
 \operatorname{Hom}_{\mathcal N}(U\cdot A,B)
 \cong\operatorname{Hom}_{\mathbb Z}
       \bigl(U,\operatorname{Hom}_{\mathcal N}(A,B)\bigr).
\]
A map from the torsion group $B$ into $\operatorname{Hom}(U,A)$ lands in its torsion subgroup, which proves this formula directly. We have natural identifications
\[
 (U\otimes V)\cdot A\cong U\cdot(V\cdot A),
 \qquad\mathbb Z\cdot A\cong A.
\]
For the first, both sides are the bilinear maps $U\times V\to A$ killed by one common positive integer. The identifications satisfy the associativity and unit coherence relations because they are the usual tensor--Hom adjunction identifications. The scalar action preserves colimits in $U$ by the copower formula. It preserves colimits in $A\in\mathcal N$ because $t\operatorname{Hom}(U,-)$ on $\mathcal T$ is right adjoint to $U\otimes -$ on $\mathcal T$.

Now put
\[
 \mathcal M=\Ch_{\geq0}(\mathbb Z)\times\mathcal N,
\]
and define a symmetric monoidal product by
\[
 (X,A)\otimes(Y,B)
 =\bigl(X\otimes_{\mathbb Z}Y,
        (X_0\cdot B)\oplus(Y_0\cdot A)\bigr).
 \tag{C-square-zero}\label{Csquarezero}
\]
Its unit is $\mathbf1=(\mathbb Z[0],0)$. To check associativity explicitly, the second coordinate of either parenthesized triple product is canonically the sum
\[
 (X_0\otimes Y_0)\cdot C\ \oplus\
 (X_0\otimes Z_0)\cdot B\ \oplus\
 (Y_0\otimes Z_0)\cdot A.
\]
The scalar-action coherence gives the pentagon; permutation of these summands and the usual signed symmetry of chain complexes give the hexagon and symmetry. The unit identities follow from $\mathbb Z\cdot A=A$ and $0\cdot A=0$. Thus \eqref{Csquarezero} is a $\mathbb Z$-bilinear symmetric monoidal tensor. It preserves colimits in both variables: tensor products of complexes, degree-zero evaluation, and the two scalar actions all do so. The category $\mathcal M$ is complete and cocomplete, abelian, and has enough projectives, since both factors do. Its products are exact, by exactness of products of abelian groups and exactness of direct sums in $\mathcal T$.

The scalar-copower functor and its two adjoints are
\[
 \kappa(U)=(U[0],0),\qquad
 \varphi(X,A)=H_0(X),\qquad
 \operatorname{Hom}_{\mathcal M}(\mathbf1,(X,A))=X_0.
\]
The functor $\kappa$ preserves all small limits as well as colimits and is strong monoidal. Projection $F:\mathcal M\to\Ch_{\geq0}(\mathbb Z)$, $(X,A)\mapsto X$, is an exact cocontinuous resolving functor for $\varphi$: on projective objects its positive homology is zero by Claim~\ref{C1}. Hence $L\varphi(X,A)$ is canonically $X$, and the normalized counit $L\varphi(\mathbf1)\to\mathbb Z$ is the identity.

Take $J=(D^1\mathbb Z,0)$ and let $\alpha:\mathbf1\to J$ be the standard degree-zero inclusion. Since $(D^1\mathbb Z)_0=\mathbb Z$, tensoring it with $(X,A)$ gives
\[
 (X,A)\longrightarrow\bigl(D^1\mathbb Z\otimes X,A\bigr),
\]
whose first coordinate is the degreewise split inclusion and whose second coordinate is $1_A$. It is monic. The first coordinate of its target is contractible, by tensoring the usual contraction of $D^1\mathbb Z$ with $X$. Thus $L_i\varphi(J\otimes(X,A))=0$ for every $i\geq0$, as required by (d).

Nevertheless $\mathcal N$, and hence $\mathcal M$, fails AB4. Fix a prime $p$ and consider the epimorphisms $q_n:\mathbb Z/p^n\to\mathbb Z/p$, $n\geq1$, in $\mathcal T$. Their product is
\[
 t\!\left(\prod_{n\geq1}\mathbb Z/p^n\right)
 \longrightarrow (\mathbb Z/p)^{\mathbb N}.
\]
The constant sequence $(1,1,\ldots)$ has no preimage. Any lift has $n$th coordinate prime to $p$ and therefore of order exactly $p^n$; such a sequence has infinite order and does not belong to the torsion subgroup. The product map is not epic. Opposite variance gives monomorphisms in $\mathcal N$ whose coproduct is not monic. In particular AB5 fails too.

Finally $F$ and its section $S(X)=(X,0)$ induce inverse equivalences after localization: $FS=1$, and the natural map $(X,A)\to(X,0)$ lies in $W$ because its image under $F$ is $1_X$. We therefore obtain $\mathcal M[W^{-1}]\simeq D_{\geq0}(\mathbb Z)$ directly. This last verification separates the failure of AB4 from the truth of the desired equivalence in this example.
\end{proof}

